\section{Choice of the height-fiber kernels}\label{sec:08-kernels}

We construct normal kernels that preserve height, make each
height-kernel image finite, and meet the transfer tests already fixed.
Their restrictions to actual incidence groups must also give the exact
homological graph covers used in the later cubulation. The construction
therefore combines finite-index separations with characters and their
nontrivial cyclic-cover cohomology modes.

Fix one periodic peripheral \(D\), its height kernel \(S=D\cap L\),
and its actual invariant subtree \(Z_D\subset T\).  The conclusions of
Sections~\ref{sec:03-rows}, \ref{sec:06-closure}, and
\ref{sec:07-transfer} used here are the following.  The group \(D\) is
torsion-free hyperbolic, quasiconvex in \(G\), and virtually compact
special; \(h(D)\ne0\).  Its action on \(Z_D\) is cocompact, and every
vertex and edge stabilizer
\[
                         M_\xi=D_\xi=S\cap X_\xi
\]
is nontrivial, finitely generated, and of type \(F\).  These stabilizers
are quasiconvex in their actual rows.  The local period rule holds at
every vertex and in either slot, and the subgroup division fields and
the injectivity of the ambient \(2\)-columns are those already
established.  In particular the argument below uses the actual rows
\(X_\xi\), not their folded replacements.

All extra requirements in this section form finite lists up to the
original group action.  They include prescribed finite avoidance,
cyclic membership and double-coset tests from
Section~\ref{sec:transfer-relative}, every actual incidence stabilizer,
any additional required nontrivial groups,
specified cyclic generators in \(S\), and pairs of actual cyclic
subgroups with trivial intersection.  For such a pair, the maximal
cyclic closures are distinct, even when they are conjugate in \(D\).
We also allow arbitrary fixed lower bounds on the girths of the
incidence core graphs and on specified cyclic orders, and prescribed
divisibility of those orders.  All these requirements are fixed before
the choices below and imposed before the normal finite-index base
subgroup \(H\unlhd D\) is frozen.  The final kernel \(N\) is then
constructed with the exact restriction identity
\eqref{kernel:local-N}, which gives the required graph covers.
Intersecting \(N\) with an arbitrary further finite-index subgroup
would preserve finite exclusions but need not preserve this identity.

\subsection{Local cones and finite-index preparations}

\begin{lemma}[The finite family of local cones]\label{kernel:local-cones}
For a vertex \(x\in Z_D\) and either slot, the incidence chain map
\[
 \bigoplus_{\substack{e\text{ in the slot at }x\\ \bmod M_x}}
       \mathbb QD\otimes_{\mathbb QM_e}C_*(M_e)
       \longrightarrow
       \mathbb QD\otimes_{\mathbb QM_x}C_*(M_x)
\]
has a finite based free mapping cone over \(\mathbb QD\) which becomes
contractible over \(D_D\).  There are finitely many such cones up to
\(D\).
\end{lemma}

\begin{proof}
Choose finite classifying complexes for the stabilizers, using their
already established torsion-free hyperbolicity and
\cite[Theorem~5]{MeintrupSchick2002}, and cellular
maps realizing their inclusions, with the actual attaching conjugations.
The orbit list at each vertex and the list of vertex types are finite.
Induction of their cellular chains therefore gives the asserted finite
free complexes with rational group-ring entries.

In degree one the induced map is an isomorphism by the local period
rule, because every incident row in the slot participates.  In degree
zero all terms vanish over \(D_D\): every stabilizer contains a
nonidentity element \(g\), and \(g-1\) is invertible in its subgroup
division field.  In degree two the homology vanishes by the induced
chain computation and injectivity of the ambient \(2\)-columns.
Higher homology vanishes by the dimension-two subgroup resolutions.
Thus the cone is acyclic over \(D_D\).  A bounded acyclic complex of
vector spaces over a division ring is contractible, by choosing
complements to the boundaries in each degree.
\end{proof}

\begin{lemma}[Novikov openness after restriction]\label{kernel:novikov-open}
There is a normal finite-index subgroup \(H_0\unlhd D\) and a nonempty
open cone
\[
                    \mathcal O_0\subset\Hom(H_0,\mathbb R)
\]
on which all the restricted local cones admit contractions over the
ordinary rational group Novikov ring.  This property persists after
passing to any further finite-index subgroup, with a possibly new
nonempty open cone.
\end{lemma}

\begin{proof}
A virtually compact special group is virtually RFRS
\cite{HaglundWise2008,Agol2008}.  Apply the finite-complex form of the
RFRS Novikov theorem \cite[Theorem~5.2]{Kielak2020} to the direct sum of
the cones in Lemma~\ref{kernel:local-cones}, after restriction to a
finite-index RFRS subgroup.  The subgroup-field compatibility gives the
required field acyclicity after restriction.  Take the normal core of
the resulting subgroup.

Here are the restriction and openness details, which also justify
later finite-index refinements.  In finite coset coordinates, a
Novikov contraction restricts to a matrix contraction over the
completion for the restricted character.  The finitely many coset
coordinate shifts do not affect summability.  Given such a contraction
\(s\), truncate its entries to group polynomials \(s_0\), far enough
that every term of
\[
                            E=1-ds_0-s_0d
\]
has strictly positive character value.  This is possible because the
differentials have finite support.  Positivity of these finitely many
terms remains true in a neighborhood of the character.  On that
neighborhood, \((1-E)^{-1}=\sum_{n\ge0}E^n\) converges in the ordinary
Novikov ring.  Since \(E\) commutes with \(d\),
\(s_0(1-E)^{-1}\) is a contraction.  Intersect the finitely many
neighborhoods for the cones and take positive radial multiples.
The same argument applies after every further restriction.  Notice
that this argument uses an actual RFRS subgroup; no selective filling
extension of the Novikov theorem is required here.
\end{proof}

\begin{lemma}[Visible and independent evaluations]\label{kernel:visibility}
After a further normal finite-index refinement of \(H_0\), the following
conditions hold, and persist under further finite-index refinement.
For each required nontrivial group \(M\), there is \(u_M\in M\cap H_0\)
on which some rational character of \(H_0\) is nonzero.  For every
specified pair \(\langle a\rangle,\langle b\rangle\) with trivial
intersection, evaluation on positive powers of \(a,b\) in \(H_0\)
defines a surjection
\[
             \Hom(H_0,\mathbb Q)\longrightarrow\mathbb Q^2.
\]
The conditions hold for all the conjugated tests as well.
\end{lemma}

\begin{proof}
Choose one nonidentity element in each required \(M\).  Its cyclic
group is quasiconvex in \(D\).  For a specified pair, the two endpoint
pairs in the boundary of \(D\) are disjoint: in a torsion-free
hyperbolic group, cyclic subgroups sharing an endpoint have a common
nontrivial power.  Their positive and negative axes consequently have
a bounded cancellation constant at a change from one generator to
the other.  Choose powers whose translation lengths exceed twice
this bound and the fixed local-to-global threshold for quasigeodesics.
The path for a reduced alternating word, after removing the bounded
pieces at its joins, is a uniform local quasigeodesic, hence a global
quasigeodesic.  No such nonempty word is trivial, and the orbit map
of the resulting rank-two free group is a quasi-isometric embedding.
Thus these powers generate a quasiconvex free group, as required.

Pass first to a compact special subgroup of finite index.
Quasiconvex subgroups of a hyperbolic compact special group are virtual
retracts \cite{HaglundWise2008,HaglundWise2012}.  Apply this to suitable
powers of each chosen cyclic group and each of the rank-two groups.
A retraction onto a cyclic group gives its nonzero rational character;
a retraction onto the rank-two group gives the two coordinate
characters.  Intersect the finitely many resulting subgroups and take
a normal core in \(D\).  Restricting these characters still detects
the required powers: a further power lying in the smaller subgroup
has its former evaluation multiplied by a nonzero integer.  The same
observation preserves independence of the two evaluation functionals.
Conjugation transports the resulting characters, proving the assertion
for all conjugates.  At a fixed normal finite-index stage only finitely
many conjugacy tests on characters are needed, since inner
conjugations by that subgroup act trivially on its characters.
\end{proof}

We may simultaneously make the finite avoidance and girth requirements
part of \(H_0\).  Indeed \(D\) is residually finite and its quasiconvex
subgroups are separable.  Products and double cosets of the cyclic
subgroups in question are separable by quasiconvex product separability
\cite{Minasyan2006}.  For a failed membership test \(u\notin Q tR\),
choose a finite quotient preserving that failure and intersect its
kernel with the current subgroup.  Tests with a full height-fiber
factor are immediate, because all the letters in the peripheral
coordinates lie in that fiber.  For each actual incidence group \(M\)
embedded in a free row, take its finite Stallings core immersed in a
rose for that row.  Choose a base vertex and retain the incidence
identification of the core fundamental group with the actual subgroup
\(M\).  This core has only finitely many nonempty closed reduced
paths of a prescribed bounded length, up to its vertices.  Transporting
each such loop to the base vertex along a chosen core path gives a
nonidentity element of \(M\le D\).  Excluding these elements, and
using normality to exclude their conjugates, makes every induced
incidence cover have the requested girth.  Any further subgroup
preserves these exclusions.

\begin{lemma}[Separating actual cyclic pairs]\label{kernel:cyclic-separation}
There is a normal finite-index subgroup \(H\unlhd D\), contained in the
prescribed \(H_0\), such that every specified pair
\(\langle a\rangle,\langle b\rangle\) has nontrivial images with trivial
intersection in \(D/H\).  Every specified generator may additionally
be required to have order divisible by a prescribed integer, and
larger than a prescribed bound.
\end{lemma}

\begin{proof}
List the conjugacy types of the maximal cyclic groups containing the
specified generators.  In a torsion-free hyperbolic group these form
a malnormal quasiconvex family.  Apply the malnormal special quotient
theorem \cite[Corollary~2.8]{AGM2016} with sufficiently deep power
kernels in these maximal cyclic groups, retaining its designated
finite-index normal subgroups $\dot P_i\unlhd P_i$, and also apply the relative
Dehn filling theorem \cite[Theorem~1.1]{Osin2007}.  We obtain a
hyperbolic virtually compact special quotient \(D'\) in which the
chosen maximal cyclic groups have the exact prescribed finite orders.
For $P_i=\langle c_i\rangle$, choose its power exponent divisible by
both $[P_i:\dot P_i]$ and $[P_i:P_i\cap H_0]$, as well as by all
prescribed order factors. Increase these multiples to meet the
finite avoidance and lower-bound requirements. The power kernels
then lie in both $\dot P_i$ and $H_0$, as required by the two
quotient constructions. Thus if \(a=c^k\), with \(c\) a maximal cyclic generator,
choosing the order of \(c\) divisible by \(|k|r\) makes the order of
\(a\) divisible by \(r\).

We verify separation in \(D'\), including when the two marked maximal
cyclic groups are conjugate.  Use the coned relative Cayley graph,
with its uniform hyperbolicity constant and the uniform
original-peripheral angle dictionaries of
Section~\ref{sec:transfer-relative}.  The constants come from fixed relative
relators and the uniform linear relative isoperimetric bound
\cite[Lemmas~2.7 and~5.3]{Osin2007}; hence they are chosen before the
power exponents.  Preserve the distinctness of the finitely many
marked apex cosets by the finite peripheral membership tests.  Let
\(v\ne w\) be two resulting marked apices, and let \(I\) be the
intersection of their stabilizers.  For \(u\in I\), compare a geodesic
\(\gamma\) from \(v\) to \(w\) with \(u\gamma\).  Thin bigons bound the
angle between their first germs by a constant independent of \(u\)
and of the filling exponents.  For a short geodesic use the path
through \(w\); otherwise join the two geodesics a fixed distance from
\(v\).  This gives an avoiding path of uniformly bounded length.

Normalize \(v\) to the apex of a maximal cyclic group
\(P=\langle c\rangle\).  Its first ordinary neighbor has coordinate
\(x\in\bar P\), so the first-neighbor displacement for \(u\gamma\) is
\(x^{-1}ux=u\), since \(\bar P\) is cyclic.  The uniform angle dictionary
therefore represents every \(u\in I\) by one of a fixed finite list of
powers \(c^e\), say \(|e|\le B\).  If \(|\bar P|=n>3B\), no nontrivial
subgroup of \(\bar P\) is contained in these residues: every nontrivial
subgroup of a cyclic group of order \(n\) has an element at cyclic
distance at least \(n/3\) from the identity.  Thus \(I=1\).
This proves the required pairwise separation in \(D'\).

Finally use residual finiteness of \(D'\).  For each of the finitely
many marked finite subgroup pairs \(A,B\le D'\), separate from the
identity every nonidentity element \(xy^{-1}\), \(x\in A,\ y\in B\);
also separate the nonidentity elements of each marked finite subgroup.
The resulting finite quotient preserves both the exact orders and
the trivial pairwise intersections.  Take its product with the
existing quotient \(D/H_0\), which factors through \(D'\) because the
power-filling kernel lies in \(H_0\).  The kernel of the resulting map
from \(D\) is the required \(H\).
\end{proof}

Freeze \(H\) after imposing all these finite requirements.  By
Lemma~\ref{kernel:novikov-open}, let
\(\mathcal O\subset\Hom(H,\mathbb R)\) be a nonempty open cone of
simultaneous contraction characters.  We may include its \(D\)-translates:
the local cone family was induced from \(D\), so conjugation preserves
the contraction property.  Write
\[
                 S_0=H\cap S,\qquad M^0=M\cap H.
\]

\subsection{Integral characters and scalar local systems}

\begin{lemma}[A finite generic character family]\label{kernel:characters}
There is a nonempty finite family of integral characters
\(\chi_j:H\to\mathbb Z\), each giving all the local Novikov
contractions, such that:
\begin{enumerate}[label=\textup{(\roman*)}]
\item every \(\chi_j\) is nonzero on each required \(M^0\) and on the
      required power in \(H\) of every specified cyclic generator;
\item for each specified pair \(a,b\), if \(m_a,m_b\) are their orders
      modulo \(H\), then
\[
       \bigl(\chi_j(a^{m_a})\bigr)_j,\qquad
       \bigl(\chi_j(b^{m_b})\bigr)_j
\]
      are linearly independent over \(\mathbb Q\);
\item the family is closed under conjugation by \(D\).
\end{enumerate}
There are arbitrarily large odd primes \(\ell\) for which all these
nonzero evaluations remain nonzero and all these pairs remain
independent modulo \(\ell\).
\end{lemma}

\begin{proof}
The forbidden zero evaluations are finitely many proper rational
hyperplanes by Lemma~\ref{kernel:visibility}.  For each pair, the
evaluation map to \(\mathbb R^2\) is onto.  Thus the condition that the
evaluations of a tuple of at least two characters have rank less than
two is a proper algebraic condition.  A nonzero real polynomial cannot
vanish on a nonempty open set.  Consequently a rational tuple in a
sufficient finite power of \(\mathcal O\) avoids all the hyperplanes
and all the rank-drop conditions simultaneously.  Scale its members
by positive integers to obtain integral characters.

Impose the same conditions on the finitely many conjugated tests,
then add all \(D\)-conjugates of the characters.  These still have the
contraction property; the orbit of an integral character is finite
because \(H\) is normal and its inner conjugations fix every character.
Adding coordinates cannot destroy the existing rank-two conditions.
For reduction modulo \(\ell\), exclude the finitely many primes
dividing any of the chosen nonzero evaluations or a chosen nonzero
\(2\)-by-\(2\) minor for each pair.
\end{proof}

\begin{lemma}[Scalar specialization]\label{kernel:scalar-specialization}
One can choose such an odd prime \(\ell\), followed by a finite field
\(\mathbb k\) of arbitrarily large characteristic different from
\(\ell\), containing all \(\ell\)-th roots of unity, so that for every
\(j\) and every \(\lambda^\ell=1\), \(\lambda\ne1\), the local
restriction maps
\[
 H^1(M_x^0;\mathbb k_{\lambda^{\chi_j}})
       \longrightarrow
 \bigoplus_{\substack{e\text{ in the chosen slot at }x\\
                         \bmod M_x^0}}
 H^1(M_e^0;\mathbb k_{\lambda^{\chi_j}})
\]
are isomorphisms.  Actual attaching conjugations are included in this
notation.  All these local groups have zero \(H^0\) for these
coefficients.
\end{lemma}

\begin{proof}
For an integral character \(\chi:H\to\mathbb Z\), there is a ring map
from the ordinary group Novikov ring to the Laurent-series field:
\[
 \widehat{\mathbb QH}^{\,\chi}\longrightarrow\mathbb Q((t)),
 \qquad
 \sum_{g\in H}r_g g\longmapsto
        \sum_{n\in\mathbb Z}
          \left(\sum_{\chi(g)=n}r_g\right)t^n .
\]
Each inner sum is finite and the occurring heights are bounded below.
The same support condition makes every coefficient of a product a
finite sum, so this is indeed a ring homomorphism.  Apply it to a
Novikov contraction of each restricted cone.  The resulting finite
complex is contractible over \(\mathbb Q((t))\).  Its differentials
already lie in \(\mathbb Q[t,t^{-1}]\); their ranks over
\(\mathbb Q(t)\) and over \(\mathbb Q((t))\) agree, as can be read from
their minors.  Hence the complex is acyclic over \(\mathbb Q(t)\).

Choose a rational-function contraction and clear all its
denominators.  For the finitely many cones and characters this gives
a nonzero Laurent polynomial \(f(t)\in\mathbb Z[t,t^{-1}]\) such that
all the contractions specialize wherever \(f(t)\ne0\).  Include
\(f(t^{-1})\) in this list as well, to accommodate the dual coefficient
convention in cohomology.  Only finitely many complex roots of unity
are roots of these polynomials.  Choose the prime \(\ell\) from
Lemma~\ref{kernel:characters} so large that no nonidentity
\(\ell\)-th root is among those roots.  Equivalently, each relevant
polynomial is relatively prime over \(\mathbb Q\) to the cyclotomic
polynomial \(\Phi_\ell(t)\).  Clearing a Bezout identity shows that
this remains true modulo every prime outside a finite set.  Exclude
also the finitely many rational coefficient denominators and
\(\ell\).  Choose any remaining sufficiently large characteristic
and a finite extension containing the \(\ell\)-th roots.  All the
specialized contractions survive in this field.

For completeness, the local maps asserted in the statement are
exactly the maps controlled by these cones.  Restricting an induced
complex from \(D\) to \(H\) decomposes it over the double cosets
\(H\backslash D/M_x\).  Within one endpoint summand, the edge terms
decompose over the incident-edge orbits of \(H\cap M_x\).
Thus each restricted cone is a finite direct sum of the actual local
cones at the corresponding \(H\)-orbit representatives.  Scalar
acyclicity gives acyclicity in every summand.  Dualizing finite
complexes over \(\mathbb k\), with inverse characters if required,
gives precisely the displayed cohomology isomorphism.
Finally each stabilizer contains an element with nonzero
\(\chi_j\)-value modulo \(\ell\).  Its scalar action
\(\lambda^{\chi_j}\) is nontrivial, since \(\ell\) is prime, so there
are no invariant vectors.
\end{proof}

\begin{proposition}[Restriction to one actual level]\label{kernel:level-cohomology}
For every \(j,\lambda\) as in
Lemma~\ref{kernel:scalar-specialization}, and every vertex or edge level
\(r\) of \(Z_D\), restriction is an isomorphism
\[
 H^1(S_0;\mathbb k_{\lambda^{\chi_j}})
       \longrightarrow
 \prod_{\substack{\xi\text{ at level }r\\\bmod S_0}}
       H^1(M_\xi^0;\mathbb k_{\lambda^{\chi_j}}).
 \tag{K1}\label{kernel:level-equation}
\]
The product on the right is finite.  In particular the group on the
left is finite-dimensional, and restriction to any individual actual
stabilizer is onto.
\end{proposition}

\begin{proof}
The \(H\)- and \(S_0\)-stabilizers of an actual position coincide,
because actual stabilizers have height zero.  Each level has finite
quotient by \(S_0\).  To see this directly, in a fixed \(D\)-orbit two
positions at the same height differ by an element of \(S\);
therefore \(D\)-cocompactness gives finitely many \(S\)-orbits per
level, and \(S_0\) has finite index in \(S\).

Take products of the local restriction isomorphisms over all vertex
orbits at one level.  On either slot the incident-edge indexing is
exactly the \(S_0\)-orbit indexing at the neighboring edge level.
The finite-index induction decomposition in
Lemma~\ref{kernel:scalar-specialization} ensures that the coefficient
transport and the attaching frames agree in every summand.
Thus restriction identifies the product of the vertex-level
cohomologies isomorphically with the product at either neighboring
edge level.

The cohomological Mayer--Vietoris sequence for the action on the tree
\(Z_D\) uses products over all vertex and edge orbits
\cite{Brown1982,Serre1980}.  All its local degree-zero terms vanish.
Its degree-one part therefore identifies
\(H^1(S_0;\mathbb k_{\lambda^{\chi_j}})\) with the kernel of
\[
 \prod_{\xi\text{ vertex}\bmod S_0}H^1(M_\xi^0;\mathbb k_{\lambda^{\chi_j}})
 \longrightarrow
 \prod_{\xi\text{ edge}\bmod S_0}H^1(M_\xi^0;\mathbb k_{\lambda^{\chi_j}}),
\]
where the map is the difference of the two endpoint restrictions.
Prescribing the data at one level determines the data at successive
levels uniquely, in both directions, by the preceding isomorphisms.
Every such prescription gives a compatible product.  This proves
\eqref{kernel:level-equation}; it also explains why infinite products
over the whole height line cause no extra degrees of freedom.
\end{proof}

\subsection{The exact kernel and its graph covers}

For every \(j\), define
\[
 S_{\ell j}=\ker\bigl(S_0\xrightarrow{\chi_j}\mathbb Z
                                \longrightarrow\mathbb Z/\ell\bigr),
 \qquad
 M_{\ell j}=M^0\cap S_{\ell j}.
\]
The selected nonzero evaluations imply that both \(S_0\) and every
required \(M^0\) map onto this same cyclic deck group.  Conjugation
therefore defines a deck action on their first cohomology with trivial
\(\mathbb k\)-coefficients.  If \(\delta\) is a deck generator, put
\[
\begin{aligned}
 e_0&=\frac1\ell\sum_{r=0}^{\ell-1}\delta^r,\\
 V_j&=\ker\bigl(e_0:H^1(S_{\ell j};\mathbb k)
                          \to H^1(S_{\ell j};\mathbb k)\bigr),\\
 V_{Mj}&=\ker\bigl(e_0:H^1(M_{\ell j};\mathbb k)
                          \to H^1(M_{\ell j};\mathbb k)\bigr).
\end{aligned}
\]
Thus these spaces contain all, and only, the nontrivial deck modes.

\begin{proposition}[Exact local annihilators]\label{kernel:exact-kernel}
Each \(V_j\) is finite-dimensional and restriction maps it onto
\(V_{Mj}\) for every actual stabilizer \(M\).  With
\(\ker V_j=\bigcap_{v\in V_j}\ker v\), the subgroup
\[
                     N=\bigcap_j\bigl(S_{\ell j}\cap\ker V_j\bigr)
                     \tag{K2}\label{kernel:N}
\]
is normal in \(D\), has finite index in \(S\), and satisfies
\[
 M\cap N
       =\bigcap_j\left(M_{\ell j}\cap
                         \bigcap_{v\in V_{Mj}}\ker v\right).
                     \tag{K3}\label{kernel:local-N}
\]
The equality uses no simultaneous surjectivity across different \(j\).
\end{proposition}

\begin{proof}
Finite-index Shapiro and the decomposition of the regular
\(\mathbb k[\mathbb Z/\ell]\)-module into its characters give
\[
 H^1(S_{\ell j};\mathbb k)
       \cong
       \bigoplus_{\lambda^\ell=1}
                    H^1(S_0;\mathbb k_{\lambda^{\chi_j}}).
\]
This remains valid when \(S_0\) is not finitely generated; the index is
finite and the decomposition has only \(\ell\) summands.
The nontrivial summands form \(V_j\), up to the harmless inverse
indexing convention for the deck characters.  They are
finite-dimensional by Proposition~\ref{kernel:level-cohomology}.

Since \(M^0\) surjects onto the full deck group, the corresponding
Shapiro decomposition for \(M_{\ell j}\) has the same character
indexing.  Naturality identifies restriction on each summand with
\[
 H^1(S_0;\mathbb k_{\lambda^{\chi_j}})
                 \longrightarrow
 H^1(M^0;\mathbb k_{\lambda^{\chi_j}}),
\]
which is onto by Proposition~\ref{kernel:level-cohomology}.
Restriction is deck-equivariant (a deck generator may be represented
inside \(M^0\)), so its image on \(V_j\) is exactly \(V_{Mj}\).

With trivial coefficients, first cohomology is
\(\Hom(S_{\ell j},\mathbb k^+)\).  A basis of \(V_j\) thus gives a
homomorphism to a finite additive group, and \(\ker V_j\) has finite
index in \(S_{\ell j}\).  There are finitely many \(j\), proving
\([S:N]<\infty\).  Conjugation by \(D\) permutes the characters,
the subgroups \(S_{\ell j}\), and their nontrivial-mode spaces.
Consequently it preserves \(N\).
Finally surjectivity \(V_j\to V_{Mj}\) says, for each \(j\) separately,
that the annihilator of the restricted global classes is the
annihilator of all local nontrivial classes.  Intersecting these
individual equalities proves \eqref{kernel:local-N}.
\end{proof}

\begin{corollary}[The outgoing graph-cover property]\label{kernel:graph-cover}
Let \(M\) be an incidence group embedded in a free actual row and
represented by its finite Stallings core.  A component \(L_H\) of the
cover determined by \(M^0\) may be required to have arbitrarily large
prescribed girth.  It has a nonempty finite family of connected cyclic
\(\ell\)-covers \(L_{\ell j}\), corresponding to \(M_{\ell j}\).
The component \(L_N\) determined by \(M\cap N\) is exactly the cover
whose fundamental group is the intersection of the covering kernels
and the kernels of every additive first-cohomology class in the
nontrivial deck part of each \(L_{\ell j}\), over the finite field
\(\mathbb k\) with \(\operatorname{char}\mathbb k\ne\ell\).

In particular, if \(e_1,e_2\) are two lifts to \(L_{\ell j}\) of the
same oriented edge of \(L_H\), the pullback to \(L_N\) of
\(e_1^*-e_2^*\) has zero period on every closed path.
All these statements are invariant under the required coset and
conjugation changes.
\end{corollary}

\begin{proof}
Connectedness of \(L_{\ell j}\) is the surjectivity of
\(\chi_j:M^0\to\mathbb Z/\ell\).  The exact description of its further
cover is \eqref{kernel:local-N}, using
\(H^1(L_{\ell j};\mathbb k)=\Hom(M_{\ell j},\mathbb k^+)\).
The graph girth was imposed at the \(H\) stage.
The deck average of \(e_1^*-e_2^*\) is zero: each lifted edge in the
deck orbit occurs once with each sign.  Since a graph has no
\(2\)-cells, this cochain is a cocycle and its class belongs to the
nontrivial deck part.  Equation~\eqref{kernel:local-N} therefore
annihilates all its periods on \(L_N\).  Transport of core graphs and
local systems preserves every step.
\end{proof}

\begin{proposition}[Orders and exact cyclic intersections]\label{kernel:orders}
For every specified generator \(a\in S\), let \(m_a\) be its order in
\(D/H\).  Its order in \(S/N\) is exactly \(m_a\ell\).  For every
specified pair of actual noncommensurate cyclics, their images in
\(S/N\) have trivial intersection.  Thus the prescribed evenness,
largeness, and divisibility of cyclic orders survive the construction.
\end{proposition}

\begin{proof}
Put \(z=a^{m_a}\in S_0\).  Every \(\chi_j(z)\) is nonzero modulo
\(\ell\), so \(z\) represents a deck generator for
\(S_0/S_{\ell j}\).  The element \(z^\ell\) lies in \(S_{\ell j}\)
and is fixed by conjugation by \(z\).  For a nontrivial deck
eigenvector \(v\in V_j\), this implies
\(v(z^\ell)=\lambda v(z^\ell)\) for some \(\lambda\ne1\), hence
\(v(z^\ell)=0\).  Decomposing \(V_j\) into eigenspaces gives
\(z^\ell\in\ker V_j\) for every \(j\), so \(a^{m_a\ell}\in N\).
Conversely \(a^q\in N\subset H\) implies \(q=m_ar\), and the condition
\(a^q\in S_{\ell j}\), for any \(j\), gives
\(r\chi_j(z)=0\bmod\ell\).  Therefore \(\ell\mid r\), proving the
exact order.

Suppose next that \(a^qN=b^rN\) for a specified pair.  The images of
their cyclic groups in \(D/H\) have trivial intersection by
Lemma~\ref{kernel:cyclic-separation}; hence \(q=m_au\) and \(r=m_bv\).
Since \(a^qb^{-r}\in S_{\ell j}\) for every \(j\),
\[
              u\bigl(\chi_j(a^{m_a})\bigr)_j
                  -v\bigl(\chi_j(b^{m_b})\bigr)_j=0
                       \quad\text{over }\mathbb F_\ell .
\]
The two vectors are independent, so \(\ell\mid u,v\).
The exact order calculation then gives \(a^q,b^r\in N\).  Their
common image was therefore trivial.
\end{proof}

\begin{remark}[Additional maximal cyclic peripherals]\label{kernel:cyclic-peripheral}
If an added peripheral is itself \(P=\langle c\rangle\) at height
zero, choose \(H_P=\langle c^r\rangle\) with \(r\) as large and as
divisible as all its finite tests require.  Use the character
\(\chi(c^r)=1\) and choose \(\ell\) avoiding the finitely many indices
\([H_P:M\cap H_P]\) for its required nontrivial incidence subgroups.
Then
\[
                              N_P=\langle c^{r\ell}\rangle .
\]
The character is surjective modulo \(\ell\) on each such incidence.
The deck action on the first cohomology of a cyclic kernel is
trivial, so its nontrivial mode part is zero; nevertheless the exact
cover description in Corollary~\ref{kernel:graph-cover} still holds.
In particular edge-lift differences can be nonzero exact cochains
on these cyclic covers.  Taking \(r\) divisible by each relevant
generator exponent and by the requested order factors gives all
the same order and girth requirements.  This applies to the cyclic
root incidences of the no-hit root expansion.
\end{remark}

\begin{proposition}[Simultaneous compatible kernels]\label{kernel:construction}
For the finite orbit list of periodic peripherals and any added
maximal cyclic peripherals, all the preceding choices can be made
simultaneously and equivariantly.  The resulting \(N_D\unlhd D\)
have finite index in \(S_D\), satisfy every prescribed finite
membership, product, avoidance, girth and cyclic-order requirement,
and give the exact local graph covers of
Corollary~\ref{kernel:graph-cover}.  Specified noncommensurate cyclic
intersections remain exact; commensurate cyclic intersections satisfy
the usual intersection identity in a finite cyclic quotient.
\end{proposition}

\begin{proof}
Make every finite requirement at the \(H_0\) and \(H\) stages,
including all its finitely many orbit representatives and
transporters.  Perform the character choices for each periodic
peripheral.  There are finitely many evaluations, minors and
specialization polynomials over the entire list, so one prime
\(\ell\) and then one finite field \(\mathbb k\) can serve all of
them, including Remark~\ref{kernel:cyclic-peripheral}.  Extend the
kernels by conjugation.  This is well-defined because each kernel
is normal in its full peripheral stabilizer.

Every final kernel lies in the corresponding frozen \(H\), so a
membership or product exclusion holding modulo \(H\) also holds
modulo that kernel: an equality in the finer quotient would project
to the forbidden equality.  The exact noncommensurate intersections
and the cyclic orders are Proposition~\ref{kernel:orders}.  For
commensurate cyclics, write their subgroups as \(m\mathbb Z,n\mathbb Z\)
in their common maximal cyclic group, and let the kernel there be
\(k\mathbb Z\).  The intersection of their images has generator
\(\operatorname{lcm}(\gcd(m,k),\gcd(n,k))\), which equals
\(\gcd(\operatorname{lcm}(m,n),k)\); prime by prime, this is the
distributive identity for minimum and maximum.  It is therefore
exactly the image of their original intersection.  Full-factor
intersection tests are immediate.  The graph-cover conclusion is
Corollary~\ref{kernel:graph-cover} and its cyclic case above.
\end{proof}
